\documentclass[11pt,a4paper]{article}

% ── Geometry & fonts ──────────────────────────────────────────────────────────
\usepackage[top=1in, bottom=1in, left=1.25in, right=1.25in]{geometry}
\usepackage[T1]{fontenc}
\usepackage[utf8]{inputenc}
\usepackage{lmodern}
\usepackage{microtype}

% ── Layout helpers ────────────────────────────────────────────────────────────
\usepackage{parskip}
\usepackage{enumitem}
\usepackage{titlesec}
\usepackage[hidelinks]{hyperref}

% ── Spacing ───────────────────────────────────────────────────────────────────
\setlength{\parindent}{0pt}
\setlength{\parskip}{0pt}

% Section headings: large bold, unnumbered, compact
\titleformat{\section}{\large\bfseries}{}{0em}{}
\titlespacing*{\section}{0pt}{14pt}{6pt}

% ─────────────────────────────────────────────────────────────────────────────
\begin{document}

% ── Title block ───────────────────────────────────────────────────────────────
\begin{center}
  {\LARGE SoS 2026: Algorithmic Trading}\\[6pt]
  {\large\textbf{Plan of Action}}\\[10pt]
  Arush Anand\\
  Mentor: Vedant Maheshwari\\[6pt]
  May 23, 2026
\end{center}

\bigskip

% ── Timeline ──────────────────────────────────────────────────────────────────
\section{Timeline}

% Reusable description-list style (Week N | topics)
\newlist{weeklist}{description}{1}
\setlist[weeklist]{
  labelwidth  = 5em,
  leftmargin  = 5.5em,
  labelsep    = 0.5em,
  style       = sameline,
  font        = \normalfont\bfseries,
  topsep      = 0pt,
  parsep      = 0pt,
  itemsep     = 5pt,
}

\begin{weeklist}

\item[Week 1] Orientation \& First Data: what quant and algorithmic trading
  actually is and who does it. Study just enough market vocabulary for a coder:
  instruments, exchanges, returns, volatility; pull real historical price data
  in Python, clean it, and plot it. \textit{(Chan Ch.~1)}

\item[Week 2] Fishing for Ideas: where strategies come from, plausible
  strategies and their pitfalls, the intuition behind mean reversion vs.
  momentum. Hand-code the signal logic for a simple moving-average crossover
  and run it on the data from Week~1. \textit{(Chan Ch.~2)}

\item[Week 3] Backtesting I: Backtesting platforms and
  data sources, writing my own simple backtester in Python, key performance
  metrics (Sharpe ratio, maximum drawdown, annualised returns); get a full
  end-to-end backtest running on my Week-2 strategy. \textit{(Chan Ch.~3)}

\item[Week 4] Backtesting II: Biases that make backtests lie
  look-ahead, survivorship, data-snooping plus transaction costs and
  slippage; deliberately expose and then fix these flaws in my Week-3
  backtest; re-evaluate performance honestly. \textit{(Chan Ch.~3)}

\end{weeklist}

\begin{center}
  \textbf{Mid-term Report Submission}\\
  \textit{\small Deliverable: a working, honestly-evaluated backtested strategy}
\end{center}

\begin{weeklist}

\item[Week 5] Mean Reversion \& Pairs Trading: build a second, more
  interesting strategy end-to-end, cointegration testing, spread
  construction, entry/exit signals, backtest and compare against Week-2
  strategy; bridges Chan's practice into Cartea's statistical-arbitrage
  framing. \textit{(Chan Ch.~2--3; Cartea pairs-trading chapter)}

\item[Week 6] Market Microstructure \& the Order Book: Study electronic market structure,
  the limit order book (LOB), order types (market, limit, cancel), bid--ask
  spread, price formation, adverse selection. \textit{(Cartea Ch.~1--2)}

\item[Week 7] Optimal Execution \& Liquidation: price impact (temporary
  vs.\ permanent), the optimal liquidation problem, Almgren--Chriss-style
  execution trajectories, trade scheduling to minimise market impact; the
  mathematics here directly answers the slippage problems encountered in
  Week~4. \textit{(Cartea Ch.~5--8)}

\item[Week 8] Market Making, Risk \& Wrap-up: inventory-based and
  adverse-selection market-making models; money management, position sizing,
  the Kelly criterion, leverage; assemble the full pipeline. \textit{(Cartea Ch.~10; Chan Ch.~6--8)}

\end{weeklist}

\begin{center}
  \textbf{End-term Report Submission}\\
  \textit{\small Deliverable: full pipeline + write-up of what you built and learned}
\end{center}

\bigskip

% ── References ────────────────────────────────────────────────────────────────
\section{References}

\textit{Algorithmic and High-Frequency Trading} by
\textit{\'Alvaro Cartea, Sebasti\'an Jaimungal \& Jos\'e Penalva}

\medskip

\textit{Quantitative Trading: How to Build Your Own Algorithmic Trading
Business} (2nd ed.)\ by \textit{Ernest P.\ Chan}

\end{document}
